\chapter{Differential Geometry for General Relativity}

\section{Purpose and Scope}

General relativity is not a theory of forces acting in a fixed Euclidean arena.
It is a theory in which spacetime itself is represented by a differentiable
manifold equipped with a Lorentzian metric. Matter and energy determine the
curvature of this metric, while freely falling bodies move along curves selected
by the metric connection. For this reason, the mathematical language of general
relativity is the language of manifolds, tensor fields, connections, geodesics,
and curvature.

This chapter gives a mathematically precise introduction to the differential
geometric structures needed before one can state Einstein's field equation in a
coordinate-independent way. The discussion is intentionally formal, but each
formal object is also connected to its physical role. The point is not to make
curvature intuitive by vague analogies. The point is to define the objects
correctly and then explain how those objects become physical in relativity.


The guiding chain of ideas is
\[
\text{smooth manifold}
\longrightarrow
\text{tangent spaces}
\longrightarrow
\text{tensor fields}
\longrightarrow
\text{metric}
\longrightarrow
\text{connection}
\longrightarrow
\text{curvature}
\longrightarrow
\text{Einstein tensor}.
\]
A reader who understands this chain understands the mathematical skeleton of
general relativity.

\begin{remarkbox}
The word \emph{manifold} should not be interpreted as a surface sitting inside a
larger Euclidean space. An abstract manifold is a space which locally has
coordinates, but it does not need an ambient space. In general relativity,
spacetime is normally treated as an abstract four-dimensional manifold.
\end{remarkbox}

\clearpage

\section{Topological and Smooth Manifolds}

A manifold is first a topological object. Smoothness is additional structure. It
is important not to confuse the two levels. Topology tells us which points are
near which other points; smooth structure tells us which coordinate changes are
differentiable.

\begin{definitionbox}
An \emph{\(n\)-dimensional topological manifold} is a topological space \(M\)
such that:
\begin{enumerate}
    \item \(M\) is Hausdorff;
    \item \(M\) is second countable;
    \item every point \(p\in M\) has an open neighbourhood \(U\subseteq M\)
    homeomorphic to an open subset of \(\mathbb{R}^n\).
\end{enumerate}
A pair \((U,\varphi)\), where \(U\subseteq M\) is open and
\(\varphi:U\to \varphi(U)\subseteq \mathbb{R}^n\) is a homeomorphism, is called
an \emph{coordinate chart}.
\end{definitionbox}

If \(p\in U\), the coordinate map gives
\[
\varphi(p)=(x^1(p),\ldots,x^n(p)).
\]
The functions \(x^1,\ldots,x^n\) are the local coordinate functions. They are not
intrinsic physical quantities. They are labels assigned to events or points by a
chosen chart.

Two charts \((U,\varphi)\) and \((V,\psi)\) overlap on \(U\cap V\). On the
overlap, the coordinate transition map is
\[
\psi\circ \varphi^{-1}:\varphi(U\cap V)\to \psi(U\cap V).
\]
This map expresses the new coordinates as functions of the old coordinates. The
ability to differentiate tensor fields depends on the differentiability of these
transition maps.

\begin{definitionbox}
A \emph{smooth atlas} on \(M\) is a collection of coordinate charts
\(\{(U_\alpha,\varphi_\alpha)\}\) covering \(M\) such that every transition map
\[
\varphi_\beta\circ \varphi_\alpha^{-1}
\]
is a \(C^\infty\) diffeomorphism between open subsets of \(\mathbb{R}^n\). A
\emph{smooth manifold} is a topological manifold together with a maximal smooth
atlas.
\end{definitionbox}

In relativity, the smooth structure allows us to write differentiable fields:
worldlines, vector fields, tensor fields, metrics, connections, and curvature.
Without this structure, the expression \(\partial g_{\mu\nu}/\partial x^\rho\)
would have no invariant meaning.

\clearpage

\section{Coordinate Independence and Diffeomorphisms}

The central mathematical discipline in differential geometry is to distinguish
objects from their coordinate representations. A physical spacetime event is not
an ordered quadruple of numbers. A coordinate system assigns such quadruples to
events, but a different coordinate system assigns different quadruples to the
same events.

\begin{definitionbox}
Let \(M\) and \(N\) be smooth manifolds. A map \(F:M\to N\) is \emph{smooth} if,
for every pair of charts \((U,\varphi)\) on \(M\) and \((V,\psi)\) on \(N\) with
\(F(U)\subseteq V\), the coordinate representation
\[
\psi\circ F\circ \varphi^{-1}
\]
is a smooth map between open subsets of Euclidean spaces.
\end{definitionbox}

A diffeomorphism is a smooth bijection with smooth inverse. If
\(F:M\to N\) is a diffeomorphism, then \(M\) and \(N\) have the same smooth
structure for all differential-geometric purposes.

In general relativity, diffeomorphisms have a special conceptual role. They
express the fact that the theory is not tied to a preferred coordinate system.
The field equations should be equations between geometric objects on a manifold,
not equations whose meaning depends on one selected coordinate grid.

\begin{remarkbox}
Coordinate covariance alone is not the same as physical insight. One can write
bad physics in a covariant notation. The serious point is that the fundamental
fields in general relativity are geometric objects: the metric, curvature,
stress-energy tensor, and the maps or curves describing matter.
\end{remarkbox}

If \(x^\mu\) and \(x^{\mu'}\) are two coordinate systems on an overlap, then a
coordinate transformation is written locally as
\[
x^{\mu'}=x^{\mu'}(x^1,\ldots,x^n).
\]
The Jacobian matrix is
\[
\frac{\partial x^{\mu'}}{\partial x^\nu}.
\]
The transformation rules for vectors, covectors, and tensors are built from this
Jacobian and its inverse. This is why tensor notation is the correct notation
for coordinate-independent physics.

\clearpage

\section{Tangent Vectors as Derivations}

In Euclidean space, a tangent vector can be drawn as an arrow. On an abstract
manifold there is no ambient vector space in which such an arrow naturally
lives. Therefore the tangent vector must be defined intrinsically.

\begin{definitionbox}
Let \(M\) be a smooth manifold and let \(p\in M\). A \emph{tangent vector at
\(p\)} is a linear map
\[
X_p:C^\infty(M)\to \mathbb{R}
\]
satisfying the Leibniz rule
\[
X_p(fh)=f(p)X_p(h)+h(p)X_p(f)
\]
for all smooth functions \(f,h\in C^\infty(M)\). The vector space of tangent
vectors at \(p\) is denoted by \(T_pM\).
\end{definitionbox}

This definition encodes directional differentiation. A tangent vector is an
operator which differentiates functions at a point.

In local coordinates \((x^1,\ldots,x^n)\), the coordinate tangent vectors are
\[
\left.\frac{\partial}{\partial x^1}\right|_p,
\ldots,
\left.\frac{\partial}{\partial x^n}\right|_p.
\]
Every tangent vector can be written uniquely as
\[
X_p=X^\mu\left.\frac{\partial}{\partial x^\mu}\right|_p,
\]
where the Einstein summation convention is used: repeated upper and lower
indices are summed.

If coordinates are changed from \(x^\mu\) to \(x^{\mu'}\), the coordinate basis
transforms by the chain rule:
\[
\frac{\partial}{\partial x^{\mu'}}
=
\frac{\partial x^\nu}{\partial x^{\mu'}}
\frac{\partial}{\partial x^\nu}.
\]
Therefore the components of a vector transform as
\[
X^{\mu'}=
\frac{\partial x^{\mu'}}{\partial x^\nu}X^\nu.
\]
The vector itself is invariant; only its coordinate components change.

\clearpage

\section{Cotangent Vectors and Differentials}

The dual space of the tangent space is the cotangent space.

\begin{definitionbox}
The \emph{cotangent space} at \(p\in M\) is
\[
T_p^*M=\operatorname{Hom}(T_pM,\mathbb{R}).
\]
Its elements are called \emph{cotangent vectors}, \emph{covectors}, or
\emph{one-forms} at \(p\).
\end{definitionbox}

If \((x^1,\ldots,x^n)\) is a coordinate system, the dual basis of \(T_p^*M\) is
\[
dx^1|_p,\ldots,dx^n|_p,
\]
where
\[
dx^\mu\left(\frac{\partial}{\partial x^\nu}\right)=\delta^\mu_{\nu}.
\]
Here \(\delta^\mu_{\nu}\) is the Kronecker delta.

For a smooth function \(f\), its differential at \(p\) is the covector
\[
df_p:T_pM\to \mathbb{R}
\]
defined by
\[
df_p(X_p)=X_p(f).
\]
In coordinates,
\[
df=\frac{\partial f}{\partial x^\mu}dx^\mu.
\]
Thus gradients in differential geometry are first naturally covectors, not
vectors. A metric is needed to convert a covector into a vector.

Under coordinate transformation, a covector \(\omega=\omega_\mu dx^\mu\) has
components transforming as
\[
\omega_{\mu'}=
\frac{\partial x^\nu}{\partial x^{\mu'}}\omega_\nu.
\]
This is the inverse transformation compared with vector components. This is why
upper and lower indices have different transformation laws.

\clearpage

\section{Tensor Spaces}

A tensor at a point is a multilinear map built from tangent and cotangent
spaces. Tensors are the natural algebraic objects of coordinate-independent
physics.

\begin{definitionbox}
A tensor of type \((r,s)\) at \(p\in M\) is an element of
\[
T^r_s(T_pM)=
\underbrace{T_pM\otimes\cdots\otimes T_pM}_{r\text{ factors}}
\otimes
\underbrace{T_p^*M\otimes\cdots\otimes T_p^*M}_{s\text{ factors}}.
\]
Equivalently, it may be viewed as a multilinear map with \(r\) covector slots
and \(s\) vector slots.
\end{definitionbox}

A type \((1,1)\) tensor can be viewed as a linear map
\[
A:T_pM\to T_pM.
\]
A type \((0,2)\) tensor is a bilinear form
\[
B:T_pM\times T_pM\to \mathbb{R}.
\]
A type \((2,0)\) tensor is a bilinear form on covectors.

In coordinates, a tensor of type \((r,s)\) has components
\[
T^{\mu_1\cdots\mu_r}{}_{\nu_1\cdots\nu_s}.
\]
Under a change of coordinates, each upper index receives a factor
\(\partial x^{\mu'}/\partial x^\mu\), and each lower index receives a factor
\(\partial x^\nu/\partial x^{\nu'}\). For example, a type \((1,1)\) tensor
transforms as
\[
A^{\mu'}{}_{\nu'}=
\frac{\partial x^{\mu'}}{\partial x^\alpha}
\frac{\partial x^\beta}{\partial x^{\nu'}}
A^\alpha{}_{\beta}.
\]
This transformation law is not a convention. It is what makes the tensor an
intrinsic object independent of coordinates.

\clearpage

\section{Tensor Fields}

A tensor field assigns a tensor smoothly to each point of a manifold. Physical
fields in relativity are tensor fields on spacetime.

\begin{definitionbox}
A \emph{tensor field of type \((r,s)\)} on a smooth manifold \(M\) is a smooth
assignment
\[
p\mapsto T_p\in T^r_s(T_pM).
\]
The vector space of smooth type \((r,s)\) tensor fields is often denoted by
\(\mathcal{T}^r_s(M)\).
\end{definitionbox}

Examples include:
\[
X=X^\mu\frac{\partial}{\partial x^\mu}
\]
for a vector field,
\[
\omega=\omega_\mu dx^\mu
\]
for a one-form, and
\[
T=T_{\mu\nu}dx^\mu\otimes dx^\nu
\]
for a covariant two-tensor field.

In relativity, the stress-energy tensor \(T_{\mu\nu}\) is a symmetric
\((0,2)\)-tensor field. The metric \(g_{\mu\nu}\) is also a symmetric
\((0,2)\)-tensor field, but with a special nondegeneracy and signature property.
Curvature tensors are built from the metric and its associated connection.

A coordinate expression such as \(T_{00}\) is not by itself a scalar physical
quantity. It is a component relative to a chosen coordinate basis. Meaningful
statements must be made either tensorially or by specifying the observer or
frame relative to which the component is measured.

\clearpage

\section{Vector Fields and Integral Curves}

A vector field may be understood as an infinitesimal flow on a manifold. At each
point it gives a direction and rate of change.

\begin{definitionbox}
A \emph{vector field} on \(M\) is a smooth section of the tangent bundle
\(TM\). In local coordinates it has the form
\[
X=X^\mu\frac{\partial}{\partial x^\mu},
\]
where the component functions \(X^\mu\) are smooth.
\end{definitionbox}

An integral curve of \(X\) is a smooth curve \(\gamma:I\to M\) satisfying
\[
\dot{\gamma}(t)=X_{\gamma(t)}.
\]
In coordinates this becomes the system of ordinary differential equations
\[
\frac{dx^\mu}{dt}=X^\mu(x^1(t),\ldots,x^n(t)).
\]
The existence and uniqueness theorem for ordinary differential equations implies
local existence and uniqueness of integral curves for smooth vector fields.

\begin{theorembox}
Let \(X\) be a smooth vector field on \(M\), and let \(p\in M\). There exists an
open interval \((-\varepsilon,\varepsilon)\) and a unique integral curve
\(\gamma\) with \(\gamma(0)=p\), defined at least for sufficiently small time.
\end{theorembox}

Physically, a vector field can represent a velocity field of a continuous
medium, a timelike observer congruence, or a generator of a symmetry. In
relativity, a timelike vector field may represent a family of observers, while a
Killing vector field represents an infinitesimal spacetime symmetry.

\clearpage

\section{Lie Brackets and Noncommuting Directions}

Vector fields act on smooth functions as derivations. Given two vector fields
\(X\) and \(Y\), their commutator measures the failure of differentiating first
along \(X\) and then along \(Y\) to equal differentiating in the opposite order.

\begin{definitionbox}
The \emph{Lie bracket} of two vector fields \(X\) and \(Y\) is the vector field
\([X,Y]\) defined by
\[
[X,Y](f)=X(Y(f))-Y(X(f))
\]
for every smooth function \(f\).
\end{definitionbox}

If
\[
X=X^\mu\partial_\mu,
\qquad
Y=Y^\mu\partial_\mu,
\]
then
\[
[X,Y]^\mu=X^\nu\partial_\nu Y^\mu-Y^\nu\partial_\nu X^\mu.
\]
The Lie bracket is bilinear, antisymmetric, and satisfies the Jacobi identity
\[
[X,[Y,Z]]+[Y,[Z,X]]+[Z,[X,Y]]=0.
\]

Coordinate vector fields commute:
\[
\left[\frac{\partial}{\partial x^\mu},
\frac{\partial}{\partial x^\nu}\right]=0.
\]
A general frame need not commute. Noncommuting frames are common in physics,
especially in orthonormal tetrad formalisms, rotating frames, and gauge-theoretic
formulations.

The Lie bracket is defined without a metric and without a connection. It belongs
to the differentiable structure of the manifold.

\clearpage

\section{Differential Forms}

Differential forms are antisymmetric covariant tensor fields. They provide the
natural language for integration on manifolds, orientation, volume forms, and
field strengths.

\begin{definitionbox}
A \emph{differential \(k\)-form} on \(M\) is a smooth section of
\(\Lambda^kT^*M\). Equivalently, it is an antisymmetric covariant tensor field
of rank \(k\).
\end{definitionbox}

A one-form is written
\[
\alpha=\alpha_\mu dx^\mu.
\]
A two-form is written
\[
\omega=\frac{1}{2}\omega_{\mu\nu}dx^\mu\wedge dx^\nu,
\]
where \(\omega_{\mu\nu}=-\omega_{\nu\mu}\). The wedge product is antisymmetric:
\[
dx^\mu\wedge dx^\nu=-dx^\nu\wedge dx^\mu.
\]

The exterior derivative is a coordinate-independent operator
\[
d:\Omega^k(M)\to \Omega^{k+1}(M)
\]
satisfying
\[
d^2=0.
\]
For a function \(f\),
\[
df=\partial_\mu f\,dx^\mu.
\]
For a one-form \(\alpha=\alpha_\mu dx^\mu\),
\[
d\alpha=\frac{1}{2}(\partial_\mu\alpha_\nu-
\partial_\nu\alpha_\mu)dx^\mu\wedge dx^\nu.
\]

In electromagnetism, the electromagnetic field strength is naturally a two-form
\(F\), and the homogeneous Maxwell equations are simply
\[
dF=0.
\]
This example shows why differential forms are not merely mathematical ornament:
they are structurally efficient descriptions of physical fields.

\clearpage

\section{Metrics: Riemannian and Lorentzian}

A metric is not merely a device for measuring distances. It is a smoothly varying
nondegenerate bilinear form on tangent spaces. Its signature determines the
causal character of vectors.

\begin{definitionbox}
A \emph{pseudo-Riemannian metric} on a smooth manifold \(M\) is a smooth
symmetric \((0,2)\)-tensor field \(g\) such that, for every \(p\in M\), the
bilinear form
\[
g_p:T_pM\times T_pM\to \mathbb{R}
\]
is nondegenerate.
\end{definitionbox}

Nondegenerate means that if
\[
g_p(X,Y)=0
\]
for all \(Y\in T_pM\), then \(X=0\). In coordinates,
\[
g=g_{\mu\nu}dx^\mu\otimes dx^\nu,
\]
and nondegeneracy means
\[
\det(g_{\mu\nu})\neq 0.
\]

A Riemannian metric is positive definite:
\[
g_p(X,X)>0
\]
for every nonzero \(X\in T_pM\). A Lorentzian metric on a four-dimensional
manifold has signature either \((-+++)\) or \((+---)\). In this chapter we use
the convention \((-+++)\).

For Lorentzian signature, a nonzero tangent vector \(X\) is:
\[
\begin{array}{ll}
\text{timelike} & \text{if } g(X,X)<0,\\
\text{null or lightlike} & \text{if } g(X,X)=0,\\
\text{spacelike} & \text{if } g(X,X)>0.
\end{array}
\]
This classification is the mathematical origin of causal structure in general
relativity.

\clearpage

\section{Line Element and Proper Time}

In coordinates, the metric is often written through the line element
\[
ds^2=g_{\mu\nu}dx^\mu dx^\nu.
\]
This expression is shorthand for the quadratic form determined by the metric.
It should not be understood as an equation between infinitesimal numbers in a
naive sense.

For a smooth curve \(\gamma:\lambda\mapsto x^\mu(\lambda)\), the tangent vector
has components
\[
\dot{x}^\mu=\frac{dx^\mu}{d\lambda}.
\]
The squared norm of the tangent is
\[
g(\dot{\gamma},\dot{\gamma})
=g_{\mu\nu}\frac{dx^\mu}{d\lambda}
\frac{dx^\nu}{d\lambda}.
\]
For a timelike curve in signature \((-+++)\), proper time satisfies
\[
d\tau^2=-\frac{1}{c^2}ds^2.
\]
If units with \(c=1\) are used, then
\[
d\tau=\sqrt{-g_{\mu\nu}dx^\mu dx^\nu}.
\]
Thus the metric determines the physical time measured by an ideal clock moving
along a timelike worldline.

For a null curve,
\[
ds^2=0.
\]
The light cone at a point consists of all null tangent directions. Therefore the
metric determines which events can be connected by light signals and which
curves can represent massive particles.

\clearpage

\section{Raising and Lowering Indices}

Because the metric is nondegenerate, it defines an isomorphism between tangent
and cotangent spaces:
\[
\flat:T_pM\to T_p^*M,
\qquad
X\mapsto X^\flat,
\]
where
\[
X^\flat(Y)=g(X,Y).
\]
In components,
\[
X_\mu=g_{\mu\nu}X^\nu.
\]
The inverse metric \(g^{-1}\) has components \(g^{\mu\nu}\) satisfying
\[
g^{\mu\alpha}g_{\alpha\nu}=\delta^\mu_\nu.
\]
It raises indices:
\[
\omega^\mu=g^{\mu\nu}\omega_\nu.
\]

This operation is not harmless notation. It uses the metric. Without a metric,
there is no canonical way to identify vectors and covectors.

For a rank-two tensor, one may form mixed components such as
\[
T^\mu{}_{\nu}=g^{\mu\alpha}T_{\alpha\nu}.
\]
In relativity this is common for the stress-energy tensor, the Ricci tensor, and
the Einstein tensor. The metric is therefore not simply another tensor field; it
also controls the algebra of index manipulation.

\clearpage

\section{Volume Element}

A metric determines a natural volume density. If \(g_{\mu\nu}\) is the matrix of
the metric in coordinates, write
\[
g=\det(g_{\mu\nu}).
\]
For a Lorentzian metric with signature \((-+++)\), this determinant is negative
in standard orientations, so the invariant volume element is written as
\[
dV_g=\sqrt{|g|}\,d^nx.
\]
In four-dimensional spacetime,
\[
dV_g=\sqrt{-g}\,d^4x
\]
when \(g<0\).

This factor is essential in generally covariant action principles. The
Einstein-Hilbert action is
\[
S_{\mathrm{EH}}=\frac{1}{16\pi G}\int_M R\sqrt{-g}\,d^4x,
\]
where \(R\) is the scalar curvature. The factor \(\sqrt{-g}\) makes the
integral independent of the coordinate system.

If coordinates change from \(x^\mu\) to \(x^{\mu'}\), the coordinate volume form
changes by the absolute value of the Jacobian determinant, while \(\sqrt{|g|}\)
changes by the inverse factor. Their product is invariant.

\clearpage

\section{Connections and Covariant Derivatives}

Ordinary partial derivatives of tensor components do not transform tensorially
under general coordinate changes. A connection corrects this failure and defines
differentiation of tensor fields along vector fields.

\begin{definitionbox}
An \emph{affine connection} on a smooth manifold \(M\) is a map
\[
\nabla:\mathfrak{X}(M)\times\mathfrak{X}(M)\to\mathfrak{X}(M),
\qquad
(X,Y)\mapsto \nabla_XY,
\]
satisfying, for smooth functions \(f\) and vector fields \(X,Y,Z\):
\begin{enumerate}
    \item \(\nabla_{fX+Y}Z=f\nabla_XZ+\nabla_YZ\);
    \item \(\nabla_X(Y+Z)=\nabla_XY+\nabla_XZ\);
    \item \(\nabla_X(fY)=X(f)Y+f\nabla_XY\).
\end{enumerate}
\end{definitionbox}

The connection is linear in the direction argument and satisfies a Leibniz rule
in the field being differentiated.

In a coordinate basis,
\[
\nabla_{\partial_\mu}\partial_\nu
=\Gamma^\rho{}_{\nu\mu}\partial_\rho.
\]
Many texts write the lower indices in the order \(\Gamma^\rho{}_{\mu\nu}\) using
\(\nabla_{\partial_\mu}\partial_\nu=\Gamma^\rho{}_{\mu\nu}\partial_\rho\). In this
chapter we use the common relativity convention
\[
\nabla_\mu V^\rho=\partial_\mu V^\rho+\Gamma^\rho{}_{\mu\nu}V^\nu.
\]
The coefficients \(\Gamma^\rho{}_{\mu\nu}\) are the Christoffel symbols of the
connection in the chosen coordinates.

\clearpage

\section{Covariant Derivative of Tensor Fields}

For a vector field \(V=V^\mu\partial_\mu\), the covariant derivative is
\[
\nabla_\mu V^\rho=\partial_\mu V^\rho+\Gamma^\rho{}_{\mu\nu}V^\nu.
\]
For a one-form \(\omega=\omega_\nu dx^\nu\), the covariant derivative is
\[
\nabla_\mu\omega_\nu=\partial_\mu\omega_\nu-
\Gamma^\rho{}_{\mu\nu}\omega_\rho.
\]
The sign is negative because covectors transform oppositely to vectors.

For a type \((1,1)\) tensor \(A^\rho{}_{\sigma}\),
\[
\nabla_\mu A^\rho{}_{\sigma}
=
\partial_\mu A^\rho{}_{\sigma}
+\Gamma^\rho{}_{\mu\lambda}A^\lambda{}_{\sigma}
-\Gamma^\lambda{}_{\mu\sigma}A^\rho{}_{\lambda}.
\]
The rule is systematic: add one connection term for each upper index and
subtract one connection term for each lower index.

The covariant derivative is constructed so that \(\nabla_\mu T^{\alpha\cdots}{}_{\beta\cdots}\)
transforms as a tensor with one additional lower index. This is the essential
reason it replaces partial derivatives in curved spacetime.

\begin{remarkbox}
The Christoffel symbols themselves are not tensor components. They can vanish at
a chosen point in suitable coordinates, but their derivatives generally cannot
all be made to vanish. Curvature is built from combinations of first derivatives
and quadratic products of Christoffel symbols that together transform tensorially.
\end{remarkbox}

\clearpage

\section{Torsion}

An affine connection has torsion if the antisymmetric part of its action on
vector fields differs from the Lie bracket.

\begin{definitionbox}
The \emph{torsion tensor} of a connection \(\nabla\) is the type \((1,2)\)
tensor field
\[
T(X,Y)=\nabla_XY-\nabla_YX-[X,Y].
\]
\end{definitionbox}

In a coordinate basis, where \([\partial_\mu,\partial_\nu]=0\), the components are
\[
T^\rho{}_{\mu\nu}=\Gamma^\rho{}_{\mu\nu}-\Gamma^\rho{}_{\nu\mu}.
\]
The connection is torsion-free precisely when
\[
\Gamma^\rho{}_{\mu\nu}=\Gamma^\rho{}_{\nu\mu}
\]
in every coordinate chart.

Standard general relativity uses the torsion-free Levi-Civita connection. Some
extensions of relativity, such as Einstein-Cartan theory, allow torsion and
relate it to spin density. In the standard theory, however, torsion is set to
zero, and the gravitational field is represented by curvature of the metric
connection.

\clearpage

\section{Metric Compatibility}

A connection is compatible with a metric when parallel transport preserves inner
products.

\begin{definitionbox}
A connection \(\nabla\) is \emph{metric-compatible} with \(g\) if
\[
\nabla g=0.
\]
In coordinates this means
\[
\nabla_\lambda g_{\mu\nu}=0.
\]
\end{definitionbox}

Expanding the coordinate expression gives
\[
\nabla_\lambda g_{\mu\nu}
=
\partial_\lambda g_{\mu\nu}
-\Gamma^\rho{}_{\lambda\mu}g_{\rho\nu}
-\Gamma^\rho{}_{\lambda\nu}g_{\mu\rho}.
\]
Thus metric compatibility is
\[
\partial_\lambda g_{\mu\nu}
=
\Gamma^\rho{}_{\lambda\mu}g_{\rho\nu}
+\Gamma^\rho{}_{\lambda\nu}g_{\mu\rho}.
\]
Geometrically, this condition means that lengths and angles, or in Lorentzian
geometry causal inner products, are preserved under parallel transport.

In relativity, metric compatibility means that the connection used to define
free fall is consistent with the metric used to define clocks, rulers, and light
cones. Without compatibility, these structures would not be tied together in the
standard way.

\clearpage

\section{The Levi-Civita Connection}

The central connection in Riemannian and Lorentzian geometry is the unique
connection that is both torsion-free and metric-compatible.

\begin{theorembox}
Let \((M,g)\) be a pseudo-Riemannian manifold. There exists a unique affine
connection \(\nabla\) such that:
\[
T=0,
\qquad
\nabla g=0.
\]
This connection is called the \emph{Levi-Civita connection} of \(g\).
\end{theorembox}

In coordinates, its Christoffel symbols are
\[
\boxed{
\Gamma^\rho{}_{\mu\nu}
=rac{1}{2}g^{\rho\sigma}
\left(
\partial_\mu g_{\nu\sigma}
+\partial_\nu g_{\mu\sigma}
-\partial_\sigma g_{\mu\nu}
\right).
}
\]

This formula follows by writing metric compatibility in three permutations of
the indices and combining them while using the symmetry
\(\Gamma^\rho{}_{\mu\nu}=\Gamma^\rho{}_{\nu\mu}\).

The Levi-Civita connection is not an additional independent field once the
metric is given. In standard general relativity, the metric determines the
connection, and the connection determines the curvature. Thus the gravitational
field is encoded in the metric.

\clearpage

\section{Parallel Transport}

Given a connection, one can define what it means for a vector to be transported
along a curve without changing relative to the connection.

Let \(\gamma:I\to M\) be a smooth curve with tangent vector \(\dot{\gamma}\). A
vector field \(V(t)\) along \(\gamma\) is parallel if
\[
\nabla_{\dot{\gamma}}V=0.
\]
In coordinates, this becomes
\[
\frac{dV^\rho}{dt}+\Gamma^\rho{}_{\mu\nu}
\frac{dx^\mu}{dt}V^\nu=0.
\]
This is a system of linear ordinary differential equations along the curve.

Parallel transport depends on the path in a curved manifold. If a vector is
parallel transported around a closed loop, it may return rotated or boosted
relative to its initial value. This failure of path independence is measured by
curvature.

In general relativity, parallel transport formalizes comparison of vectors at
different spacetime points. Since tangent spaces at different points are
different vector spaces, there is no canonical comparison without additional
structure. The connection supplies such a comparison.

\clearpage

\section{Geodesics}

A geodesic is a curve whose tangent vector is parallel transported along itself.
It is the curved-space generalization of a straight line with constant velocity.

\begin{definitionbox}
A curve \(\gamma:I\to M\) is an \emph{affinely parametrized geodesic} of a
connection \(\nabla\) if
\[
\nabla_{\dot{\gamma}}\dot{\gamma}=0.
\]
\end{definitionbox}

In coordinates, this equation is
\[
\boxed{
\frac{d^2x^\rho}{d\lambda^2}
+\Gamma^\rho{}_{\mu\nu}
\frac{dx^\mu}{d\lambda}
\frac{dx^\nu}{d\lambda}=0.
}
\]

For the Levi-Civita connection of a Lorentzian metric, timelike geodesics are
the worldlines of freely falling massive test particles, while null geodesics
are the paths of ideal light rays in geometric optics.

The word \emph{test} is important. A test particle is assumed not to disturb the
metric appreciably. The geodesic equation describes motion in a given geometry;
the Einstein equation describes how matter determines that geometry.

\clearpage

\section{Geodesics from a Variational Principle}

For a Riemannian metric, geodesics locally extremize length. In Lorentzian
geometry, timelike geodesics locally extremize proper time under suitable
conditions.

Consider the energy functional
\[
E[\gamma]=\frac{1}{2}\int_{\lambda_0}^{\lambda_1}
g_{\mu\nu}(x)
\frac{dx^\mu}{d\lambda}
\frac{dx^\nu}{d\lambda}\,d\lambda.
\]
The Euler-Lagrange equations for the Lagrangian
\[
L=\frac{1}{2}g_{\mu\nu}(x)\dot{x}^\mu\dot{x}^\nu
\]
are
\[
\frac{d}{d\lambda}\left(g_{\rho\nu}\dot{x}^\nu\right)
-rac{1}{2}\partial_\rho g_{\mu\nu}\dot{x}^\mu\dot{x}^\nu=0.
\]
Expanding the derivative gives
\[
g_{\rho\nu}\ddot{x}^\nu
+\partial_\sigma g_{\rho\nu}\dot{x}^\sigma\dot{x}^\nu
-rac{1}{2}\partial_\rho g_{\mu\nu}\dot{x}^\mu\dot{x}^\nu=0.
\]
Multiplying by \(g^{\rho\alpha}\) and symmetrizing the terms in
\(\dot{x}^\mu\dot{x}^\nu\) yields
\[
\ddot{x}^\alpha+\Gamma^\alpha{}_{\mu\nu}
\dot{x}^\mu\dot{x}^\nu=0.
\]
Thus the variational and connection definitions agree for the Levi-Civita
connection.

\clearpage

\section{Normal Coordinates and the Equivalence Principle}

At any point of a pseudo-Riemannian manifold, one can choose coordinates in
which the metric takes its flat canonical form and the Christoffel symbols of
the Levi-Civita connection vanish at that point.

\begin{theorembox}
Let \((M,g)\) be a pseudo-Riemannian manifold and let \(p\in M\). There exist
local coordinates around \(p\) such that
\[
g_{\mu\nu}(p)=\eta_{\mu\nu},
\qquad
\Gamma^\rho{}_{\mu\nu}(p)=0,
\]
where \(\eta_{\mu\nu}\) is the diagonal flat metric with the same signature as
\(g\).
\end{theorembox}

These are normal coordinates at \(p\). In Lorentzian geometry they are closely
related to local inertial frames.

However, the first derivatives of the connection cannot generally be made to
vanish at the same point. The curvature tensor remains. This distinction is
physically crucial: gravitational acceleration can be transformed away at a
point by going to a freely falling frame, but tidal effects cannot be eliminated
if curvature is nonzero.

This is the precise mathematical content behind the local equivalence principle:
locally at a point, freely falling coordinates remove the connection
coefficients, but they do not remove curvature.

\clearpage

\section{Curvature as Noncommutation of Covariant Derivatives}

Curvature measures the failure of second covariant derivatives to commute, after
correcting for the noncommutation of the vector fields themselves.

\begin{definitionbox}
The \emph{Riemann curvature tensor} of a connection \(\nabla\) is the map
\[
R(X,Y)Z=
\nabla_X\nabla_YZ-
\nabla_Y\nabla_XZ-
\nabla_{[X,Y]}Z.
\]
\end{definitionbox}

For each pair \((X,Y)\), the map \(Z\mapsto R(X,Y)Z\) is a linear endomorphism
of the tangent space. In coordinates,
\[
R(\partial_\mu,\partial_\nu)\partial_\sigma
=R^\rho{}_{\sigma\mu\nu}\partial_\rho.
\]
With the sign convention used here,
\[
\boxed{
R^\rho{}_{\sigma\mu\nu}
=
\partial_\mu\Gamma^\rho{}_{\nu\sigma}
-
\partial_\nu\Gamma^\rho{}_{\mu\sigma}
+\Gamma^\rho{}_{\mu\lambda}\Gamma^\lambda{}_{\nu\sigma}
-\Gamma^\rho{}_{\nu\lambda}\Gamma^\lambda{}_{\mu\sigma}.
}
\]

Different books sometimes use the opposite sign convention for the Riemann
tensor. What matters is internal consistency. The definitions of Ricci tensor,
scalar curvature, and Einstein tensor must follow the chosen convention.

\clearpage

\section{Curvature and Parallel Transport Around a Loop}

Curvature has a direct geometric interpretation through parallel transport. If a
vector is parallel transported around a very small coordinate parallelogram
spanned by coordinate displacements \(\Delta x^\mu\) and \(\Delta x^\nu\), the
change in the vector is, to leading order,
\[
\Delta V^\rho
\approx
R^\rho{}_{\sigma\mu\nu}V^\sigma
\Delta x^\mu\Delta x^\nu.
\]
Thus curvature measures the infinitesimal holonomy of the connection.

In flat Euclidean space, parallel transport is path-independent. On a curved
surface or curved spacetime, path dependence is the rule. The Riemann tensor is
the local object that encodes this path dependence.

Physically, this is connected to tidal effects. A uniform gravitational field
can be removed locally by choosing an accelerated or freely falling frame, but
the relative acceleration of nearby freely falling particles is controlled by
curvature and cannot be removed by a coordinate transformation.

\clearpage

\section{Algebraic Symmetries of the Riemann Tensor}

For the Levi-Civita connection of a metric, it is useful to lower the first
index:
\[
R_{\rho\sigma\mu\nu}=g_{\rho\lambda}R^\lambda{}_{\sigma\mu\nu}.
\]
The Riemann tensor then satisfies the algebraic symmetries
\[
R_{\rho\sigma\mu\nu}=-R_{\sigma\rho\mu\nu},
\]
\[
R_{\rho\sigma\mu\nu}=-R_{\rho\sigma\nu\mu},
\]
\[
R_{\rho\sigma\mu\nu}=R_{\mu\nu\rho\sigma},
\]
and the first Bianchi identity
\[
R_{\rho[\sigma\mu\nu]}=0.
\]
The brackets denote antisymmetrization over the enclosed indices.

These identities drastically reduce the number of independent components. In
\(n\) dimensions, the Riemann tensor of a Levi-Civita connection has
\[
\frac{n^2(n^2-1)}{12}
\]
independent components. In four dimensions this gives
\[
\frac{16(15)}{12}=20.
\]
Thus four-dimensional spacetime curvature has twenty independent components at
a point before field equations impose additional restrictions.

\clearpage

\section{Ricci Tensor and Scalar Curvature}

The Ricci tensor is a contraction of the Riemann tensor. With the convention of
this chapter, define
\[
\operatorname{Ric}_{\sigma\nu}=R^\rho{}_{\sigma\rho\nu}.
\]
In index notation this is often written as
\[
R_{\mu\nu}=R^\rho{}_{\mu\rho\nu}.
\]
The scalar curvature is the trace of the Ricci tensor with the inverse metric:
\[
R=g^{\mu\nu}R_{\mu\nu}.
\]
The same letter \(R\) is commonly used both for the Riemann curvature operator
and for the scalar curvature; context and indices determine the meaning.

The Ricci tensor measures part of curvature related to volume distortion and
geodesic convergence. The full Riemann tensor contains more information than the
Ricci tensor. In four-dimensional general relativity, the vacuum equation
\(R_{\mu\nu}=0\) does not imply that the full Riemann tensor is zero. Vacuum
spacetimes may still contain gravitational waves or tidal curvature.

The scalar curvature is a single number at each point. It is too compressed to
describe all curvature, but it is central in the Einstein-Hilbert action.

\clearpage

\section{Sectional Curvature and Riemannian Interpretation}

In Riemannian geometry, sectional curvature gives the curvature of a two-plane
inside a tangent space. Let \(u,v\in T_pM\) be linearly independent. The
sectional curvature of the plane they span is
\[
K(u,v)=
\frac{g(R(u,v)v,u)}{g(u,u)g(v,v)-g(u,v)^2}.
\]
This expression generalizes Gaussian curvature of surfaces.

For a two-dimensional Riemannian manifold, all curvature information is encoded
by Gaussian curvature. For dimensions three and higher, sectional curvatures for
all two-planes encode the Riemann tensor.

In Lorentzian geometry, one must be more careful because the denominator may
vanish for degenerate two-planes involving null directions. Nevertheless, the
conceptual point survives: curvature is not merely a scalar quantity. It depends
on directions and planes in tangent space.

In physics this matters because tidal effects depend on directions. Curvature
can stretch in one spatial direction and compress in another. The Riemann tensor
keeps this directional information.

\clearpage

\section{Geodesic Deviation}

Geodesic deviation gives the most direct physical interpretation of spacetime
curvature. Consider a one-parameter family of geodesics with tangent vector
\(U\) and separation vector \(\xi\). If \(U\) is affinely parametrized, the
relative acceleration satisfies
\[
\boxed{
\nabla_U\nabla_U\xi=R(U,\xi)U.
}
\]
In components, for a timelike geodesic parametrized by proper time \(\tau\),
\[
\frac{D^2\xi^\mu}{d\tau^2}
=R^\mu{}_{\nu\rho\sigma}U^\nu\xi^\rho U^\sigma,
\]
up to the sign convention for \(R\).

This equation states that curvature controls the relative acceleration of nearby
freely falling particles. This is not a coordinate artifact. It is an invariant
measurement principle: if a cloud of freely falling particles changes shape due
to tidal effects, spacetime curvature is being detected.

In a local inertial frame at a point, the Christoffel symbols can be set to zero
at that point, but the second relative acceleration encoded by geodesic
deviation remains if the Riemann tensor is nonzero.

\clearpage

\section{The Bianchi Identities}

The first Bianchi identity is algebraic:
\[
R^\rho{}_{[\sigma\mu\nu]}=0.
\]
For the Levi-Civita connection there is also a differential identity, the second
Bianchi identity:
\[
\nabla_\lambda R^\rho{}_{\sigma\mu\nu}
+
\nabla_\mu R^\rho{}_{\sigma\nu\lambda}
+
\nabla_\nu R^\rho{}_{\sigma\lambda\mu}=0.
\]
Contracting the second Bianchi identity yields
\[
\nabla^\mu\left(R_{\mu\nu}-\frac{1}{2}Rg_{\mu\nu}\right)=0.
\]
This contracted identity is one of the structural reasons why the Einstein
tensor appears in the field equations.

The Bianchi identities are not optional extra equations. They are identities
satisfied by curvature constructed from the Levi-Civita connection. Their
physical consequence is connected to local conservation of stress-energy in
general relativity.

\clearpage

\section{The Einstein Tensor}

The Einstein tensor is defined by
\[
\boxed{
G_{\mu\nu}=R_{\mu\nu}-\frac{1}{2}Rg_{\mu\nu}.
}
\]
It is symmetric because the Ricci tensor and metric are symmetric. By the
contracted Bianchi identity,
\[
\nabla^\mu G_{\mu\nu}=0.
\]
This divergence-free property is essential for coupling geometry to matter.

Einstein's field equation, with cosmological constant, is
\[
\boxed{
G_{\mu\nu}+\Lambda g_{\mu\nu}
=rac{8\pi G}{c^4}T_{\mu\nu}.
}
\]
Here \(T_{\mu\nu}\) is the stress-energy tensor. The left-hand side is geometry;
the right-hand side is matter and energy.

The equation is not simply a poetic statement that matter curves spacetime. It
is a precise nonlinear system of partial differential equations for the metric
components \(g_{\mu\nu}\), because the Ricci tensor contains second derivatives
of the metric and nonlinear terms in its first derivatives.

\clearpage

\section{Stress-Energy Tensor and Conservation}

The stress-energy tensor represents local energy density, momentum density,
energy flux, and stress as measured relative to observers or frames. In a local
orthonormal frame, components have the schematic interpretation
\[
T_{00}=\text{energy density},
\qquad
T_{0i}=\text{momentum or energy flux},
\qquad
T_{ij}=\text{stress components}.
\]
These interpretations require a frame. The tensor itself is coordinate
independent.

The Einstein equation implies
\[
\nabla^\mu T_{\mu\nu}=0
\]
provided \(\Lambda\) is constant. This is the covariant local conservation law.
It is not the same as ordinary global conservation of total energy in arbitrary
curved spacetimes. Global conservation laws require additional structure, such
as asymptotic flatness or spacetime symmetries.

For a perfect fluid with energy density \(\rho\), pressure \(p\), and four-
velocity \(u^\mu\), the stress-energy tensor is
\[
T_{\mu\nu}=(\rho+p)u_\mu u_\nu+p g_{\mu\nu}
\]
in units where \(c=1\), with signature \((-+++)\) and normalization
\[
g_{\mu\nu}u^\mu u^\nu=-1.
\]
This formula shows how matter models enter the geometric field equation.

\clearpage

\section{Local Frames and Tetrads}

A coordinate basis \(\partial_\mu\) is not usually orthonormal. In relativity it
is often useful to introduce a local frame \(e_a\), where Latin indices label
frame components. A tetrad in four dimensions is a basis
\[
e_0,e_1,e_2,e_3
\]
of each tangent space such that
\[
g(e_a,e_b)=\eta_{ab}.
\]
Here \(\eta_{ab}=\operatorname{diag}(-1,1,1,1)\) in the chosen convention.

The relation between coordinate and frame components is written
\[
e_a=e_a{}^\mu\partial_\mu.
\]
The metric can be reconstructed from tetrads by
\[
g_{\mu\nu}=\eta_{ab}e^a{}_{\mu}e^b{}_{\nu}.
\]

Physically, an orthonormal tetrad can represent the measuring frame of a local
observer: one timelike unit vector and three spacelike unit vectors. Tensor
components in this frame correspond more directly to measured quantities than
arbitrary coordinate components.

Tetrads are also important in coupling spinor fields to gravity, because spinors
are naturally associated with local Lorentz frames rather than directly with
coordinate frames.

\clearpage

\section{Killing Vector Fields and Symmetries}

A vector field generates an infinitesimal symmetry of the metric if the metric
does not change along its flow. This is expressed by the Lie derivative.

\begin{definitionbox}
A vector field \(K\) is a \emph{Killing vector field} of a metric \(g\) if
\[
\mathcal{L}_K g=0.
\]
Equivalently, for the Levi-Civita connection,
\[
\nabla_\mu K_\nu+\nabla_\nu K_\mu=0.
\]
\end{definitionbox}

Killing vector fields encode continuous spacetime symmetries. A timelike Killing
field corresponds to stationarity; a rotational Killing field corresponds to
rotational symmetry.

Along a geodesic with tangent \(U\), a Killing field gives a conserved quantity:
\[
\frac{d}{d\lambda}g(K,U)=0.
\]
This follows from the geodesic equation and the Killing equation. In physical
terms, symmetries generate conserved quantities such as energy and angular
momentum, when the relevant symmetry exists.

\clearpage

\section{Flatness and Curved Coordinates}

It is important to distinguish nonzero Christoffel symbols from genuine
curvature. Curvilinear coordinates in flat space can produce nonzero Christoffel
symbols even though the Riemann tensor is zero.

For example, the Euclidean plane in polar coordinates has metric
\[
ds^2=dr^2+r^2d\theta^2.
\]
The Christoffel symbols include
\[
\Gamma^r{}_{\theta\theta}=-r,
\qquad
\Gamma^\theta{}_{r\theta}=\Gamma^\theta{}_{\theta r}=\frac{1}{r}.
\]
These do not indicate intrinsic curvature. The Riemann tensor still vanishes.

In general relativity this distinction is essential. A gravitational field is
not identified with Christoffel symbols alone, because they can be made to vanish
at a point by choosing local inertial coordinates. Curvature, not the connection
coefficients by themselves, encodes invariant gravitational tidal effects.

\clearpage

\section{Flat Spacetime as a Lorentzian Manifold}

Special relativity is described by Minkowski spacetime. As a manifold it is
\(\mathbb{R}^4\), equipped with the Lorentzian metric
\[
\eta=-c^2dt^2+dx^2+dy^2+dz^2.
\]
In units \(c=1\),
\[
ds^2=-dt^2+dx^2+dy^2+dz^2.
\]
In standard inertial coordinates, the metric components are constant, so the
Christoffel symbols vanish:
\[
\Gamma^\rho{}_{\mu\nu}=0.
\]
Consequently,
\[
R^\rho{}_{\sigma\mu\nu}=0.
\]

However, if one uses accelerated or curvilinear coordinates on Minkowski
spacetime, Christoffel symbols may become nonzero. The curvature remains zero.
Thus special relativity is not the absence of all connection coefficients in all
coordinates; it is the vanishing of the Riemann curvature tensor of the
Minkowski metric.

\clearpage

\section{Curved Spacetime and Gravity}

In general relativity, the gravitational field is represented by a Lorentzian
metric \(g\) on a four-dimensional smooth manifold \(M\). The metric determines:
\begin{enumerate}
    \item the causal cones;
    \item proper time along timelike curves;
    \item spatial distances relative to observers;
    \item the Levi-Civita connection;
    \item the geodesics of free particles and light rays;
    \item the curvature tensors entering the field equations.
\end{enumerate}

The geometry is dynamical: the metric is not fixed in advance. It is solved from
Einstein's equation together with matter equations and appropriate initial or
boundary conditions.

The statement that gravity is curvature should be read carefully. The connection
controls free-fall equations, while curvature controls invariant tidal effects.
The metric determines both. Standard general relativity packages these facts in
the geometry of a Lorentzian manifold.

\clearpage

\section{Initial Data and the Role of Hypersurfaces}

A spacetime can often be studied by splitting it into space plus time. A
spacelike hypersurface \(\Sigma\subset M\) has an induced Riemannian metric
\(h\) and extrinsic curvature \(K\). The pair \((h,K)\) forms the geometric part
of initial data.

The Einstein equation imposes constraint equations on this data. In vacuum,
these are schematically
\[
{}^{(3)}R+K^2-K_{ij}K^{ij}=0,
\]
and
\[
D_j(K^{ij}-h^{ij}K)=0,
\]
where \({}^{(3)}R\) is the scalar curvature of the spatial metric \(h\), and
\(D\) is its Levi-Civita connection.

This shows that general relativity is not merely a set of four-dimensional
formulae. It also has an initial value formulation: one may specify suitable
data on a spatial hypersurface and evolve them, subject to constraints.

This section only indicates the structure. A rigorous treatment requires global
hyperbolicity, Sobolev spaces, and the analysis of nonlinear hyperbolic systems.

\clearpage

\section{Summary of the Mathematical Pipeline}

The mathematical construction used in general relativity can be summarized as
follows.

\begin{enumerate}
    \item A smooth four-dimensional manifold \(M\) represents spacetime as a set
    of events with differentiable structure.
    \item Tangent spaces \(T_pM\) contain possible infinitesimal directions at
    each event.
    \item Tensor fields provide coordinate-independent physical fields.
    \item A Lorentzian metric \(g\) classifies directions as timelike, null, or
    spacelike and determines proper time.
    \item The metric determines a unique Levi-Civita connection \(\nabla\).
    \item The connection defines parallel transport and geodesics.
    \item The Riemann tensor measures curvature as noncommutation of covariant
    derivatives and path dependence of parallel transport.
    \item The Ricci tensor and scalar curvature are contractions of the Riemann
    tensor.
    \item The Einstein tensor is the divergence-free curvature tensor that
    couples to stress-energy.
\end{enumerate}

\begin{summarybox}
A spacetime model in standard general relativity is a pair \((M,g)\), where
\(M\) is a smooth four-dimensional manifold and \(g\) is a Lorentzian metric.
The Levi-Civita connection of \(g\) defines free fall; its curvature defines
invariant tidal gravity; contractions of that curvature form the Einstein tensor
\[
G_{\mu\nu}=R_{\mu\nu}-\frac{1}{2}Rg_{\mu\nu},
\]
which appears in the field equation
\[
G_{\mu\nu}+\Lambda g_{\mu\nu}=\frac{8\pi G}{c^4}T_{\mu\nu}.
\]
\end{summarybox}
